\documentclass[11pt,a4paper]{article}
\usepackage[T1]{fontenc}
\usepackage[utf8]{inputenc}
\usepackage{lmodern}
\usepackage{amsmath,amssymb}
\usepackage[margin=25mm]{geometry}
\usepackage{longtable,booktabs,array,calc}
\usepackage[hidelinks]{hyperref}
\hypersetup{pdftitle={Quantum-to-Newton results: literature and Lean verification targets},pdfauthor={navstokgap research working notes}}
\newcommand{\tightlist}{\setlength{\itemsep}{0pt}\setlength{\parskip}{0pt}}
\setlength{\emergencystretch}{3em}
\setcounter{secnumdepth}{0}
\begin{document}
\hypertarget{quantum-to-newton-results-literature-and-lean-verification-targets}{%
\section{Quantum-to-Newton results: literature and Lean verification
targets}\label{quantum-to-newton-results-literature-and-lean-verification-targets}}

\textbf{Result and verification specification, 2026-10-27.} The recent
0+1 construction yields several precise propositions worth formal
checking, but their being new to this repository is not evidence of
publication priority. The strongest distinctive candidates are the
specified \(C^{1,1}\) product/packet package and its sharp-rate example,
the explicitly rescaled wall stresses, and the Gaussian preparation test
at a bounded \(C^1\) cusp. Each has substantial close prior art. The
universal constructive-lifting conjecture and an independent
positive-action necessity theorem remain \textbf{open}.

The targets below are exact mathematical specifications, \textbf{not
Lean declarations that have been compiled}. Their written proofs are in
the linked notes and are unrefereed; the \(C^{1,1}\) lift received the
internal Sol check recorded there. The targets state real quantifiers
and boundary/domain conventions so a later Lean development can
translate them without strengthening conclusions silently. A conditional
theorem whose hypotheses include the conclusion, a placeholder quantum
theory, or an unchecked imported axiom does not verify the corresponding
concrete target.

\hypertarget{what-is-already-in-the-literature}{%
\subsection{1. What is already in the
literature}\label{what-is-already-in-the-literature}}

\begin{longtable}[]{@{}
  >{\raggedright\arraybackslash}p{(\columnwidth - 4\tabcolsep) * \real{0.3333}}
  >{\raggedright\arraybackslash}p{(\columnwidth - 4\tabcolsep) * \real{0.3333}}
  >{\raggedright\arraybackslash}p{(\columnwidth - 4\tabcolsep) * \real{0.3333}}@{}}
\toprule\noalign{}
\begin{minipage}[b]{\linewidth}\raggedright
Repository proposition
\end{minipage} & \begin{minipage}[b]{\linewidth}\raggedright
Closest checked precedent
\end{minipage} & \begin{minipage}[b]{\linewidth}\raggedright
Proper status
\end{minipage} \\
\midrule\noalign{}
\endhead
\bottomrule\noalign{}
\endlastfoot
Coherent packets follow classical flow; smooth-observable semiclassics &
\href{https://doi.org/10.1007/BF01646348}{Hepp (1974)}, abstract and
theorem statements in the maintained note & Established mechanism, not a
new proof of classical mechanics as such. \\
Uniform observable time-splitting estimates &
\href{https://arxiv.org/pdf/2208.07957}{Borns--Weil--Fang, Theorem 2 and
§3.2} (passages read);
\href{https://arxiv.org/html/2507.02783v1}{Fang--Qu, Theorem 1}
(statement read) & Prior uniform-in-\(\hbar\) observable estimates. The
former uses smooth bounded symbol classes; the latter treats periodic,
spatially discretized polynomial observables and higher-order
splittings. Neither passage alone establishes the exact \(C^{1,1}\),
arbitrary-partition, coherent-packet package below. \\
Reflected Gaussian and wall impulse &
\href{https://arxiv.org/pdf/quant-ph/0408182}{Belloni--Doncheski--Robinett
(2004), equations (8)--(11), (16)--(19)} (passages read);
\href{https://www.scielo.org.mx/scielo.php?pid=S1870-35422013000200001\&script=sci_arttext}{De
Vincenzo (2013), equations (38)--(39)} (passage read) & The odd-image
Gaussian and collision-force analysis already exist. The particular
measure limit, scaled profile and Neumann comparison in Target D are
precise specializations whose exact priority remains unverified. \\
Preparation-dependent singular-flow selection &
\href{https://arxiv.org/pdf/1011.1651}{Athanassoulis--Paul, Theorem 2.1
and Lemma 5.3} (statements and relevant passages read) include cutoff
\(-|x|^{1+\theta}\), including \(\theta=1/2\);
\href{https://arxiv.org/pdf/1608.03716}{Chabu, Theorem 1.12 and §3}
(statements and passages read) gives identical incoming and different
outgoing limits for \(-|x|\). & Neither selection failure nor this
power-law cusp is new as such. Target E specifies pure Gaussian
preparations, a common incoming trajectory and their explicit mixtures
for the bounded \(C^1\) cusp. In Chabu's theorem the two profiles may
also be chosen equal, so equality of initial position densities is not
by itself a new general phenomenon. Exact coverage of Target E by prior
results has not been settled. \\
Tangent groupoid, its continuous field and strict quantization &
\href{https://arxiv.org/abs/math/9802102}{Cariñena et al.~(1998/99)},
including Rivero, and
\href{https://arxiv.org/abs/math-ph/0001005}{Landsman--Ramazan (2000)}
(abstracts and groupoid passages read) & Established background; Target
B is a particular state-level estimate on that field, not the invention
of the groupoid or of strict quantization. \\
Normalized Euclidean tower & Brownian-bridge/Feynman--Kac and martingale
constructions; \href{https://arxiv.org/abs/1108.1335}{Dimock I} and
\href{https://arxiv.org/abs/1212.5562}{II} (selected sections in the
guide) & Target G packages standard ingredients in a 0+1 illustration.
Dimock's papers treat three-dimensional scalar \(\phi^4\), not this
quantum-mechanical theorem or a four-dimensional gauge construction. \\
Relativistic one-particle limit &
\href{https://doi.org/10.1016/0022-1236(90)90049-Q}{Carmona--Masters--Simon}
and
\href{https://arxiv.org/abs/1109.5829}{Hiroshima--Ichinose--Lőrinczi}
(abstracts read) & Known relativistic Schrödinger and
subordinate-process framework. Target F is a particular joint estimate;
no literature priority has been established. \\
\end{longtable}

The table is a bounded search, not an exhaustive novelty determination.
The primary proofs and their retained limitations remain in the
\href{lipschitz-certificate-quantum-lift.md}{Lipschitz lift},
\href{quantum-blocks-construct-newton.md}{smooth construction},
\href{groupoid-quantum-newton-states.md}{groupoid note},
\href{dirichlet-wall-impulse.md}{wall note},
\href{cusp-quantum-preparation-selection.md}{cusp note}, and
\href{relativistic-dispersion-blocks.md}{relativistic note}.

\hypertarget{common-definitions-for-targets-ac}{%
\subsection{2. Common definitions for Targets
A--C}\label{common-definitions-for-targets-ac}}

Use \(\hbar>0\), \(m,\omega>0\), \(\mathcal H=L^2(\mathbb R,dq)\),
\(Qf(q)=qf(q)\), \(P_\hbar=-i\hbar\partial_q\), and
\(N_\hbar=1+Q^2+P_\hbar^2\) as the closure of the oscillator operator on
Schwartz space, with \(D(N_\hbar)=H^2(\mathbb R)\cap D(Q^2)\). Put \[
 \|(q,p)\|_\omega=|q|+|p|/(m\omega),\quad
 K_s=e^{-isV(Q)/\hbar},\quad
 D_s=e^{-isP_\hbar^2/(2m\hbar)},\quad S_s=D_sK_s.
\] For a finite positive partition \(\pi=(0=t_0<\cdots<t_n=T)\), write
\(s_j=t_{j+1}-t_j\), \(|\pi|=\max_j s_j\), and
\(S_{\hbar,\pi}(t_j)=S_{s_{j-1}}\cdots S_{s_0}\), with the empty product
at \(j=0\). The initial normalized Gaussian is \[
 g_\hbar(q)=\left(\frac{m\omega}{\pi\hbar}\right)^{1/4}
 \exp\!\left[-\frac{m\omega(q-q_0)^2}{2\hbar}
             +\frac{ip_0(q-q_0)}{\hbar}\right]. \tag{1}
\] The classical kick--drift vertices, starting at \(z_0=(q_0,p_0)\),
obey \[p^{\rm cl}_{j+1}=p^{\rm cl}_j-s_jV'(q^{\rm cl}_j),\qquad
q^{\rm cl}_{j+1}=q^{\rm cl}_j+s_jp^{\rm cl}_{j+1}/m. \tag{2}\] When used
below, \(z_{\hbar,\pi}(t)\) is the piecewise-linear interpolation of the
\textbf{quantum means at the vertices}; it is not asserted to be the
expectation of a physical intermediate state.

For the Weyl targets, use the Fourier convention \[
 \widehat b(\xi,\eta)=\int_{\mathbb R^2}
       e^{-i(\xi q+\eta p)}b(q,p)\,dq\,dp,\quad
 \operatorname{Op}^W_\hbar(b)=\frac1{(2\pi)^2}
       \int\widehat b(\xi,\eta)e^{i(\xi Q+\eta P_\hbar)}d\xi d\eta.
 \tag{3}
\] This fixes the signs and factors needed by a formalization.

\hypertarget{target-a-lipschitz-certificate-quantum-product-and-packet}{%
\subsection{3. Target A: Lipschitz certificate, quantum product, and
packet}\label{target-a-lipschitz-certificate-quantum-product-and-packet}}

\textbf{Quantifiers and assumptions.} Let
\(V\in C^1(\mathbb R;\mathbb R)\) and \(L\ge0\) satisfy
\(|V'(x)-V'(y)|\le L|x-y|\) for every \(x,y\). There is no lower bound
assumption on \(V\). For every \(m,\omega,\hbar>0\), \(T>0\) and
\(z_0\in\mathbb R^2\), the following are to be proved.

\begin{enumerate}
\def\labelenumi{\arabic{enumi}.}
\tightlist
\item
  The operator \(H_\hbar=P_\hbar^2/(2m)+V(Q)\) with initial domain
  \(\mathcal S(\mathbb R)\) is essentially self-adjoint. Its closure
  generates \(U_\hbar(t)=e^{-it\overline H_\hbar/\hbar}\), and this
  group preserves \(D(N_\hbar)\) and is strongly continuous there in the
  \(N_\hbar\) graph norm. On Schwartz space \([Q,P_\hbar]=i\hbar I\);
  for all \(u,v\in\mathbb R\) the bounded Weyl operators satisfy
  \(e^{iuQ}e^{ivP_\hbar}=e^{-i\hbar uv}e^{ivP_\hbar}e^{iuQ}\).
\item
  There exists \(0<C_{\hbar,T}<\infty\), depending on the supplied model
  and fixed \(\hbar,T\), such that \textbf{every} positive partition
  satisfies \[
   \|S_{\hbar,\pi}(T)f-U_\hbar(T)f\|_2
     \le C_{\hbar,T}|\pi|\,\|N_\hbar f\|_2
     \quad(f\in D(N_\hbar)). \tag{A1}
  \] Consequently the product converges strongly on all of
  \(\mathcal H\) as \(|\pi|\downarrow0\) at each fixed \(\hbar\). For a
  general \(\psi\in\mathcal H\) one precise modulus is the infimum over
  \(f\in D(N_\hbar)\) of
  \(2\|\psi-f\|_2+C_{\hbar,T}\delta\|N_\hbar f\|_2\). The constant in
  (A1) is \textbf{not} asserted uniform in \(\hbar\).
\item
  Let \(z_j^\hbar=(\langle Q\rangle_j,\langle P_\hbar\rangle_j)\) in the
  state \(S_{\hbar,\pi}(t_j)g_\hbar\), and let
  \(W_j=\operatorname{sd}_jQ+\operatorname{sd}_jP_\hbar/(m\omega)\). Set
  \(C=\sqrt{2/(m\omega)}\) and \(\beta=\omega+L/(m\omega)\). Then for
  every vertex, \[
   W_j\le C\sqrt\hbar\,e^{\beta t_j},\qquad
   \|z_j^\hbar-z_j^{\rm cl}\|_\omega
     \le \frac{LC}{m\omega}\,t_j e^{\beta t_j}\sqrt\hbar. \tag{A2}
  \] The classical Lipschitz flow \(\Phi_t(z_0)\) exists uniquely, and
  there is a finite classical constant \(R_T^{\rm cl}\) such that
  \(\sup_{0\le t\le T}\|z_\pi^{\rm cl}(t)-\Phi_t(z_0)\|_\omega \le R_T^{\rm cl}|\pi|\).
  Hence the vertices converge to the flow for \textbf{every} joint limit
  \(\hbar,|\pi|\downarrow0\).
\end{enumerate}

For the conjecture's bounded instance, the input is the \emph{global}
action/force and the classical polygon certificate. (A1) is proved from
that model data plus the operator toolkit; it is not a consequence of
the numerical error constant of a single orbit. Formalizing (A1) while
assuming its self-adjointness or graph bounds should be reported as a
conditional sublemma, not as the whole lift.

\hypertarget{target-b-the-smooth-ohbarpi-state-estimate}{%
\subsection{\texorpdfstring{4. Target B: the smooth \(O(\hbar+|\pi|)\)
state
estimate}{4. Target B: the smooth O(\textbackslash hbar+\textbar\textbackslash pi\textbar) state estimate}}\label{target-b-the-smooth-ohbarpi-state-estimate}}

Assume in addition that \(V\ge0\), \(V\in C^\infty\) and all its
derivatives of order at least two are bounded. Fix \(h_0,T>0\). Let
\(0<\hbar\le h_0\) and let \(\pi\) be any positive partition of
\([0,T]\). In the notation above, write \[
 \begin{gathered}
 a=\max\!\left(\omega,\frac L{m\omega}\right),\qquad
 G_T=(e^{aT}-1)/a,\\
 b_0=\frac{|V'(0)|+LC\sqrt{h_0}e^{\beta T}}{m\omega},\quad
 R_T=e^{\beta T}(\|z_0\|_\omega+Tb_0),\\
 F_T=|V'(0)|+L(R_T+C\sqrt{h_0}e^{\beta T}),\quad
 C_\delta=(F_T+LR_T)/m,\quad
 C_h=\frac{\|V'''\|_\infty C^2e^{2\beta T}}{2m\omega}.
 \end{gathered} \tag{B0}
\] The specification is \[
 \sup_{0\le t\le T}\|z_{\hbar,\pi}(t)-\Phi_t(z_0)\|_\omega
 \le G_T(C_\delta|\pi|+C_h\hbar). \tag{B1}
\] The centering step responsible for the better rate is the separate
lemma \[
 |\langle V'(Q)\rangle-V'(\langle Q\rangle)|
 \le \tfrac12\|V'''\|_\infty\operatorname{Var}(Q). \tag{B2}
\] It requires a normalized state in the needed second-moment domain.
The first derivative term cancels; the width itself remains
\(O(\sqrt\hbar)\).

For every \(b\in C_c^\infty(\mathbb R^2)\) set \[
 I_b=(2\pi)^{-2}\!\int|\widehat b|(|\xi|+m\omega|\eta|),\quad
 J_b=\frac1{2(2\pi)^2}\!\int|\widehat b|
                         (|\xi|+m\omega|\eta|)^2. \tag{B3}
\] At every partition vertex the desired \textbf{concrete} result is \[
 \left|\langle g_\hbar,S_{\hbar,\pi}(t_j)^*
       \operatorname{Op}^W_\hbar(b)S_{\hbar,\pi}(t_j)g_\hbar\rangle
       -b(\Phi_{t_j}(z_0))\right|
 \le J_b C^2e^{2\beta T}\hbar+
      I_bG_T(C_\delta|\pi|+C_h\hbar). \tag{B4}
\] It follows for any joint schedule \(\hbar,|\pi|\to0\). An additional
continuous-field corollary takes \(\hbar=\epsilon S_*\), \(S_*>0\), and
extends the limit to continuous sections of the established tangent
groupoid field: for every continuous section \(A\), every family of
partitions with \(\epsilon,|\pi_\epsilon|\to0\), and vertices
\(t_{j(\epsilon)}\to t\in[0,T]\), the corresponding vector-state value
on \(A(\epsilon)\) converges to \(A(0)(\Phi_t(z_0))\). This corollary
requires the continuous-field theorem and the identification of (3) with
the chosen quantization section; (B4) alone does not construct that
\(C^*\)-field.

\hypertarget{target-c-sharpness-of-the-square-root-rate}{%
\subsection{5. Target C: sharpness of the square-root
rate}\label{target-c-sharpness-of-the-square-root-rate}}

Set \(m=\omega=1\), \(V(q)=q|q|/2\), \(q_0=p_0=0\), and let \(g_\hbar\)
be (1). This potential is \(C^{1,1}\) and \textbf{unbounded below}. The
classical trajectory is identically \((0,0)\). Define
\[f(t)=\langle U_1(t)g_1,P_1U_1(t)g_1\rangle.\] Prove, for every
\(\hbar>0\) and real \(t\), \[
 \langle U_\hbar(t)g_\hbar,P_\hbar U_\hbar(t)g_\hbar\rangle
    =\sqrt\hbar\,f(t),\qquad f'(0)=-1/\sqrt\pi. \tag{C1}
\] In particular, there exists \(t_0>0\) such that for all
\(0<t\le t_0\) and every \(\hbar>0\), the absolute momentum error is at
least \(t\sqrt\hbar/(2\sqrt\pi)\). This refutes a uniform \(O(\hbar)\)
mean bound over the full class in Target A. It does not by itself
establish sharpness under an additional lower-bound assumption on \(V\).

\hypertarget{target-d-dirichlet-and-neumann-wall-stresses}{%
\subsection{6. Target D: Dirichlet and Neumann wall
stresses}\label{target-d-dirichlet-and-neumann-wall-stresses}}

Let \(m,\omega,q_0>0\), \(p_0<0\), \(v=p_0/m\), \(t_*=-q_0/v\),
\(A(t)=1+\omega^2t^2\), and \(\Lambda=|v|\sqrt{m\omega/A(t_*)}\). Let
\(g_{\hbar,t}\) be the free whole-line evolution of (1), and let \(R\)
reflect \(q\mapsto-q\). Put \[
 \gamma_\hbar=e^{-[m\omega q_0^2+p_0^2/(m\omega)]/\hbar},\quad
 \psi_{\hbar,D/N}(q,t)=
 \frac{g_{\hbar,t}(q)\mp g_{\hbar,t}(-q)}
      {\sqrt{1\mp\gamma_\hbar}}\quad(q>0). \tag{D0}
\] The upper sign is Dirichlet, the lower Neumann. Define the real
half-line current form and full boundary stress \[
 J_\hbar(t)=\hbar\operatorname{Im}\!\int_0^\infty
                 \overline\psi\,\partial_q\psi\,dq,\quad
 F_\hbar(t)=\frac{\hbar^2}{2m}
  \left(|\psi'(0,t)|^2-
        \operatorname{Re}[\overline{\psi(0,t)}\psi''(0,t)]\right).
 \tag{D1}
\] For both indicated wave families, prove \(J_\hbar'=F_\hbar\) and, for
every bounded continuous \(\varphi:\mathbb R\to\mathbb R\), \[
 \lim_{\hbar\downarrow0}\int_{\mathbb R}\varphi(t)F_{\hbar,D/N}(t)\,dt
    =2|p_0|\varphi(t_*). \tag{D2}
\] The Neumann measures are finite \textbf{signed} measures; their
Dirichlet-only derivative term vanishes. At each fixed \(s\in\mathbb R\)
also prove \[
 \lim_{\hbar\downarrow0}\sqrt\hbar\,
 F_{\hbar,D/N}(t_*+\sqrt\hbar\,s)
 =\frac{2p_0^2}{m}\sqrt{\frac{m\omega}{\pi A(t_*)}}
   e^{-\Lambda^2s^2},\quad
 J_{\hbar,D/N}(t_*+\sqrt\hbar\,s)\to
 |p_0|\operatorname{erf}(\Lambda s). \tag{D3}
\] This is weak convergence of measures, not convergence in total
variation to the delta. \(J_\hbar\) is a current form: the half-line
operator \(-i\hbar\partial_q\) is not being asserted self-adjoint. No
general Robin packet limit is part of this target.

\hypertarget{target-e-two-preparations-and-mixtures-at-a-c1-cusp}{%
\subsection{\texorpdfstring{7. Target E: two preparations and mixtures
at a \(C^1\)
cusp}{7. Target E: two preparations and mixtures at a C\^{}1 cusp}}\label{target-e-two-preparations-and-mixtures-at-a-c1-cusp}}

In units \(m=1\), fix \(R>0\) and a smooth even cutoff \(\chi\) equal to
one on \([-R,R]\) and zero outside \([-2R,2R]\). Set
\[V(q)=-\tfrac23|q|^{3/2}\chi(q),\quad F=-V',\quad
 H_\hbar=-\tfrac{\hbar^2}{2}\partial_q^2+V(Q),\quad
 F_*=\|F\|_\infty.\] Choose \(T>0\) such that \(F_*T^2/2<R/4\). Choose
\(C_0>0\) with \(|F(x)|\le C_0\sqrt{|x|}\) for all \(x\), and a fixed
\(K>0\) such that \(K^2>\max\{16\sqrt2C_0,16/3\}\). For
\(0<\hbar<e^{-1}\) put \(a_\hbar=1/\log(1/\hbar)\) and \[
 g_{\hbar,k}(q)=(\pi\hbar)^{-1/4}
 \exp\!\left[-\frac{(q-a_\hbar)^2}{2\hbar}
             -\frac{ik a_\hbar^{3/4}(q-a_\hbar)}{\hbar}\right],
 \quad k\in\{0,K\}. \tag{E0}
\] These two states have exactly the same initial position density. Let
\(\psi_{\hbar,k}(t)=e^{-itH_\hbar/\hbar}g_{\hbar,k}\) and set \[
 z_+(t)=(t^4/144,t^3/36),\qquad
 z_\times(t)=
 \begin{cases}z_+(t),&t\le0,\\-z_+(t),&t\ge0.\end{cases} \tag{E1}
\] The target is uniform on \([-T,T]\): \[
 (\langle Q\rangle_{\hbar,0},\langle P_\hbar\rangle_{\hbar,0})\to z_+,
 \quad
 (\langle Q\rangle_{\hbar,K},\langle P_\hbar\rangle_{\hbar,K})\to z_\times,
 \quad
 \operatorname{sd}Q_{\hbar,k}+\operatorname{sd}P_{\hbar,k}\to0.
 \tag{E2}
\] For every \(\lambda\in[0,1]\) and \(\theta\in\mathbb R\), normalize
\(\sqrt\lambda\,g_{\hbar,0}  +e^{i\theta}\sqrt{1-\lambda}\,g_{\hbar,K}\).
For every \(b\in C_c^\infty(\mathbb R^2)\), its evolved Weyl expectation
must converge uniformly on \([-T,T]\) to
\[\lambda b(z_+(t))+(1-\lambda)b(z_\times(t)). \tag{E3}\] Thus both
limiting incoming histories agree for \(t\le0\) and can have different
outgoing measures. The target neither asserts finite-\(\hbar\) Wigner
positivity nor classifies waiting-time branches.

\hypertarget{target-f-a-relativistic-joint-limit}{%
\subsection{8. Target F: a relativistic joint
limit}\label{target-f-a-relativistic-joint-limit}}

Let \(V\in C^\infty(\mathbb R;\mathbb R)\) be bounded below with all
derivatives of order at least two bounded, \(m,c_0,h_0,T>0\), and
\(z_0\in\mathbb R^2\). For \(c\ge c_0\), define
\[T_c(p)=\sqrt{m^2c^4+c^2p^2}-mc^2,\quad
 S_{\hbar,c,s}=e^{-isT_c(P_\hbar)/\hbar}e^{-isV(Q)/\hbar}. \tag{F0}\]
For the coherent packet (1), its piecewise-linear interpolated means,
and the nonrelativistic Newton flow \(\Phi_t(z_0)\), prove that there
exists \(C_T<\infty\), independent of \(\hbar,c,\pi\), such that \[
 \sup_{0\le t\le T}
 \|z_{\hbar,\pi,c}(t)-\Phi_t(z_0)\|_\omega
 \le C_T(\hbar+|\pi|+c^{-2}) \tag{F1}
\] for all \(0<\hbar\le h_0\), \(c\ge c_0\) and positive partitions of
\([0,T]\). At fixed \(\hbar\), the self-adjoint \(T_c(P_\hbar)+V(Q)\)
converges in strong resolvent sense to \(P_\hbar^2/(2m)+V(Q)\) as
\(c\to\infty\), and its unitary groups converge strongly, uniformly on
compact time intervals.

The separate free wall target uses the \emph{spectral} Dirichlet
Hamiltonian \(T_c(\sqrt{-\hbar^2\Delta_D})\), not the generally
different killed nonlocal process. For the odd packet (D0), its position
expectation satisfies, for \(|t|\le T\), \[
 \left|\langle Q\rangle^D_{\hbar,c}(t)-
           |q_0+p_0t/m|\right|
 \le\sqrt{\frac2{1-\gamma_\hbar}}
  \left[\sqrt{\frac{\hbar}{2m\omega}}+
        \frac Tm\sqrt{\frac{\hbar m\omega}{2}}\right]
       +\frac{T|p_0|^3}{2m^3c^2}. \tag{F2}
\] The positive-energy one-particle operator in an external potential is
not a construction of a Lorentz-covariant interacting field theory.

\hypertarget{target-g-an-exact-normalized-euclidean-tower}{%
\subsection{9. Target G: an exact normalized Euclidean
tower}\label{target-g-an-exact-normalized-euclidean-tower}}

Let \(V:\mathbb R\to[0,\infty)\) be continuous, \(m,\hbar,T>0\), and
\(x,y\in\mathbb R\). Let \(\mathbb P_0\) be the Brownian bridge of
variance parameter \(\hbar/m\) from \(x\) to \(y\) on \([0,T]\). Define
\[
 W=\int_0^TV(B_t)\,dt,\quad
 Z=\mathbb E_0e^{-W/\hbar},\quad
 \rho=Z^{-1}e^{-W/\hbar},\quad
 \rho_n=\mathbb E_0[\rho\mid\mathcal F_n], \tag{G0}
\] where \(\mathcal F_n\) is generated by dyadic evaluations of \(B\).
The target is \(Z>0\) and, almost surely, \[
 \mathbb E_0\rho_n=1,\quad 0<\rho_n\le Z^{-1},\quad
 \rho_n=\mathbb E_0[\rho_{n+1}\mid\mathcal F_n],\quad
 \rho_n\to\rho\ \text{a.s. and in }L^1(\mathbb P_0). \tag{G1}
\] For arbitrary positive partitions with mesh tending to zero, the
normalized densities obtained from the trapezoidal sums of \(V(B_t)\)
also converge to \(\rho\) in \(L^1(\mathbb P_0)\); their polygonal
pushforwards converge weakly on \(C([0,T])\). The pushforward measures
are not claimed to converge in total variation. The \(\rho_n\) are
conditional expectations of the continuum density; they are not
identified with the original bare finite-lattice density.

\hypertarget{the-bridge-to-newtonlean}{%
\subsection{10. The bridge to
newtonlean}\label{the-bridge-to-newtonlean}}

The inspected newtonlean master commit is
\href{https://github.com/arivero/newtonlean/commit/8138e9449b4f961c4adbebcb6090acc0d2fd5aaf}{8138e944}.
Its toolchain is Lean 4.34.1 core/Std only, with no mathlib. The current
session did not compile it. Its
\href{https://github.com/arivero/newtonlean/blob/8138e9449b4f961c4adbebcb6090acc0d2fd5aaf/BarrowLib/Polygon/FiniteEstimates.lean}{FiniteEstimates.cell}
uses \textbf{drift then kick}:
\[\operatorname{DK}_s(q,v)=(q+sv,\ v+sa(q+sv)), \tag{H0}\] whereas (2),
after \(v=p/m\) and \(a=-V'/m\), is \textbf{kick then drift}:
\[\operatorname{KD}_s(q,v)=(q+s[v+sa(q)],\ v+sa(q)). \tag{H1}\] Equality
of the two cells is false. The following finite rational lemma is an
exact interface target, requiring no quantum foundations.

\textbf{Target H.} Let \(a:\mathbb Q\to\mathbb Q\) be
\(\Lambda\)-Lipschitz, with \(\Lambda\in\mathbb Q_{\ge0}\). Let
\(s_0,\ldots,s_{n-1}\in\mathbb Q_{\ge0}\), \(\tau\in\mathbb Q_{>0}\),
and put \(\|(q,v)\|_\tau=|q|+\tau|v|\). Iterate KD and DK from the same
rational state. Suppose the DK states \((\bar q_j,\bar v_j)\) satisfy
\(|a(\bar q_j)|\le A\) and \(|\bar v_j|\le B\) for \(j<n\), with
\(A,B\in\mathbb Q_{\ge0}\). Then \[
 \|z_n^{KD}-z_n^{DK}\|_\tau
 \le(A+\tau\Lambda B)
 \sum_{j=0}^{n-1}s_j^2
  \prod_{k=j+1}^{n-1}(1+s_k/\tau)(1+\tau\Lambda s_k). \tag{H2}
\] Indeed one cell differs by at most \((A+\tau\Lambda B)s_j^2\), and KD
amplifies a previous discrepancy by at most
\((1+s_j/\tau)(1+\tau\Lambda s_j)\). The empty product is one. The
analogous statement for real data with \(\sum_js_j=T\) implies
\[\|z_n^{KD}-z_n^{DK}\|_\tau
 \le(A+\tau\Lambda B)T|\pi|e^{(1/\tau+\tau\Lambda)T}. \tag{H3}\] The
rational statement is intended as a finite algebraic lemma; (H3) also
uses real exponential bounds and is a separate completion step.

The existing
\href{https://github.com/arivero/newtonlean/blob/8138e9449b4f961c4adbebcb6090acc0d2fd5aaf/ModernLib/Polygon/GeneralForcePolygonCurve.lean}{GeneralForcePolygonCurve.polygonMap\_uniform\_convergence}
constructs a uniformly convergent position polygon on its own binary
time and Cauchy-quotient representations, under a regional
CentralOracle, force-sampling, confinement and short-window conditions.
It is not already the general one-dimensional global \(C^{1,1}\) result
in Target A: arbitrary partitions, external real-time identification and
general noncentral forces remain separate. To connect the models one
must additionally instantiate a compatible sample oracle, identify
completed coordinates and velocities with real \(q,p/m\), prove that
(H2)'s bounds hold, and show the two limits agree. The axis-embedded
harmonic oscillator is a suitable special case; general \(V\) need not
define an inward central force.

The
\href{https://github.com/arivero/newtonlean/blob/8138e9449b4f961c4adbebcb6090acc0d2fd5aaf/Reverse/Classical/ClassicalTrajectory.lean}{Reverse/Classical/ClassicalTrajectory.lean}
interfaces assume localized trajectories, free dynamics, the
momentum--velocity relation and additive impulses; the
\href{https://github.com/arivero/newtonlean/blob/8138e9449b4f961c4adbebcb6090acc0d2fd5aaf/Reverse/Parent/QuantumKKFamily.lean}{QuantumTheory}
carrier has no Hilbert-space or operator structure. Targets A--F could
provide evidence for restricted, explicitly embedded versions of some of
these interfaces. They do not automatically discharge the interfaces for
arbitrary composite bodies or prove a quantum gravity parent. Historical
Newton propositions keep their given-trajectory premise; none of these
modern targets is a historical premise.

For context only, current mathlib documentation has
\href{https://leanprover-community.github.io/mathlib4_docs/Mathlib/Analysis/ODE/ExistUnique.html}{Picard--Lindelöf},
\href{https://leanprover-community.github.io/mathlib4_docs/Mathlib/Analysis/Distribution/SchwartzSpace/Fourier.html}{Schwartz
Fourier analysis}, and
\href{https://leanprover-community.github.io/mathlib4_docs/Mathlib/Probability/BrownianMotion/Basic.html}{Brownian-motion
definitions}. Those pages are not a claim that this project has, or may
import, the needed version or all operator-theoretic results. Under the
present core-only rule, a full quantum development must construct its
analytic foundations or remain explicitly conditional; a dependency
expansion requires separate user direction. In either route, check the
actual compiled theorem types and axiom dependencies at the pinned
toolchain.

\hypertarget{consequence-for-state}{%
\subsection{11. Consequence for STATE}\label{consequence-for-state}}

These targets distinguish standard inputs, repository-specific
statements and unknown publication priority. They expose a finite
classical comparison that can be attempted under newtonlean's present
rules, and the substantial analytic foundations needed for an end-to-end
quantum verification. The universal lifting conjecture, positive-action
necessity and the four-dimensional gauge construction remain open.
\end{document}
